\documentclass[../main.tex]{subfiles}
\begin{document}

\onlyinsubfile{
  \setcounter{chapter}{1}
  \setcounter{section}{1}
  \setcounter{subsection}{0}
  \addcontentsline{toc}{chapter}{Основы кинематики}
  \addcontentsline{toc}{section}{Относительность движения}
}






%==============================================================================
\subsection{Закон сложения скоростей}
%==========================================================

Описание движения любых тел всегда ведётся относительно некоторой \textit{системы отсчёта}. Напомним, что системой отсчёта называется декартова система координат, связанная с телом отсчёта, и часы. Выбор системы отсчёта, вообще говоря, произволен: её можно связать с Землёй, с вращающимся колесом или даже с падающим камнем. Поэтому кинематические характеристики (например, траектория или скорость) одной и той же материальной точки в разных системах отсчёта обычно различаются. Впрочем это не означает, что какая-либо система отсчёта является <<более правильной>>, чем другая: все системы отсчёта равноправны. Различие состоит лишь в том, что в одной системе отсчёта движение может описываться проще и нагляднее, чем в другой. Поэтому умение переходить от одной системы отсчёта к другой является одним из основных инструментов кинематики.

Рассмотрим две системы отсчёта: лабораторную $S$ (её будем условно считаем «неподвижной») и подвижную $S'$.
Пусть В каждый момент времени положение начала координат $O'$ системы $S'$ в системе $S$ задаётся радиус-вектором $\mbfR_0(t)$. \mbox{Если} система $S'$ не вращается относительно $S$, то с точностью до поворота, можно считать их соответствующие оси сонаправленными. Положение произвольной точки $P$ характеризуется в этих системах отсчёта в момент времени $t$ радиус-векторами $\mbfr(t)$  (в системе $S$), и $\mbfr'(t)$ (в системе $S'$). Эти векторы связаны простым соотношением
\begin{equation}
\mbfr(t) = \mbfR_0(t) + \mbfr'(t),
\label{eq:r_trans_base}
\end{equation}
Пусть за малый промежуток времени от $t$ до $t + \Delta t$ векторы из формулы (\ref{eq:r_trans_base}) изменяются на: $\Delta \mbfr$, $\Delta \mbfR_0$,  $\Delta \mbfr'$ соответственно. При этом из (\ref{eq:r_trans_base}) следует равенство для малых изменений 
\begin{equation*}
\Delta \mbfr = \Delta \mbfR_0 + \Delta \mbfr',
\end{equation*}


Продифференцируем это равенство по времени:
\begin{equation*}
\dv{\mbfr}{t} = \dv{\mbfR_0}{t} + \dv{\mbfr'}{t} \implies \mbfv = \mbfv_0  + \mbfv'.
\end{equation*}
Здесь $\mbfv_0 = \dv{\vec{R}_0}{t}$~--- скорость начала координат системы $S'$ относительно $S$ (скорость переноса или \textit{переносная} скорость),
$\mbfv'$~--- скорость точки $P$ в системе $S'$
(\textit{относительная}  скорость),
$\mbfv$~--- скорость точки $P$ в системе $S$
(\textit{абсолютная}  скорость).

\begin{definition}
Закон сложения скоростей при поступательном движении одной системы отсчёта относительно другой:
\begin{equation*}
    \mbfv = \mbfv_0  + \mbfv',
\end{equation*}
где $\mbfv_0$~--- скорость системы $S'$ относительно $S$ (переносная скорость). Этот закон справедлив при любом, в том числе неравномерном, движении системы $S'$, при условии, что она движется \textit{поступательно} (без вращения).
\end{definition}

Заметим важное следствие: дифференцируя ещё раз, получаем
$\mbfa = \mbfa_0 + \mbfa'$, где $\mbfa_0 = \dv[2]{\mbfR_0}{t}$~---
ускорение системы $S'$ относительно $S$.
Если $\mbfv_0 = \mathrm{const}$ (система $S'$ движется равномерно),
то $\mbfa_0 = 0$ и ускорения совпадают в обеих системах:
$\mbfa = \mbfa'$.
Это означает, что \textit{сила, как причина ускорения, воспринимается
одинаково во всех инерциальных системах отсчёта}~---
фундаментальный принцип Галилея.

\noindent\textbf{Векторное сложение.}
Закон $\mbfv = \mbfv' + \mbfv_0$ требует именно
\textit{векторного} сложения. Решая конкретные задачи,
следует строить векторный треугольник скоростей, а не
складывать скорости алгебраически.

\begin{example}
Спортсмены бегут колонной длины $L$ со скоростью $v$.
Тренер движется им навстречу со скоростью $u$ ($u < v$),
затем, поравнявшись с последним спортсменом, разворачивается
и бежит обратно. Каждый спортсмен, поравнявшись с тренером,
разворачивается и бежит в противоположном направлении с той же
скоростью $v$. Определите длину колонны в момент, когда
тренер добежит до её головы.

Удобнее всего задачу решать в системе отсчёта, связанной
с первоначальным направлением движения колонны. В этой системе
колонна изначально покоится, её длина равна $L$. Тренер
движется к хвосту со скоростью $u + v$. Время до встречи
с хвостом:
\[
    t_1 = \frac{L}{u + v}.
\]
После встречи тренер разворачивается и движется к голове
со скоростью $v - u$ (относительно выбранной системы).
Хвост колонны, развернувшись, теперь движется в обратную
сторону со скоростью $2v$ относительно этой системы,
а голова остаётся неподвижной до момента встречи с тренером.

Время $t_2$, за которое тренер проходит расстояние от хвоста
до головы, определяется из условия, что он должен преодолеть
оставшееся расстояние, пока хвост удаляется. Однако существует
более элегантный способ: заметим, что отношение скоростей
тренера относительно колонны при движении навстречу и попутно
равно $\frac{u+v}{v-u}$. Поскольку процесс разворота каждого
спортсмена происходит последовательно и линейно во времени,
финальная длина колонны $L'$ масштабируется ровно в это же
отношение, но в обратную сторону (так как колонна «сжимается»
в процессе):
\[
    L' = L \cdot \frac{v - u}{v + u}.
\]
Таким образом, колонна укорачивается, и её новая длина определяется
соотношением относительных скоростей тренера и спортсменов.
\end{example}

\begin{example}
Бульдозер движется по горизонтальному полю со скоростью $v$.
Определите скорости относительно Земли: верхней части гусеницы,
нижней части гусеницы и точки гусеницы, скорость которой
относительно бульдозера направлена вертикально вверх.

В системе отсчёта бульдозера гусеница движется: нижняя ветвь~---
назад со скоростью $v$, верхняя~--- вперёд со скоростью $v$.
Применяем закон сложения скоростей с $\mbfv_0 = v\hat{\mbfe}_x$
(скорость бульдозера):

Нижняя ветвь: $v_{\text{ниж}} = v_{\text{бул}} - v = 0$~--- нижняя
часть гусеницы мгновенно покоится относительно Земли. Это
согласуется с тем, что нижняя часть гусеницы~--- мгновенный центр
скоростей гусеницы как тела, катящегося по Земле.

Верхняя ветвь: $v_{\text{верх}} = v + v = 2v$~--- вперёд.

Точка с вертикальной скоростью: в системе бульдозера вектор
скорости направлен вертикально вверх ($|\mbfv'| = v$).
Складываем с горизонтальной скоростью $v$ бульдозера векторно:
$|\mbfv| = \sqrt{v^2 + v^2} = v\sqrt{2}$, направление~--- под $45^\circ$
к горизонту.
\end{example}

%==============================================================================
\subsection{Метод подвижной системы отсчёта}
%==============================================================================

Многие задачи на относительное движение формально сводятся
к нахождению скорости одного тела относительно другого.
Перейдя в систему отсчёта одного из тел, мы превращаем
задачу о двух движущихся объектах в задачу об одном движущемся
и одном покоящемся. Это часто существенно упрощает геометрию.

\noindent\textbf{Минимальное расстояние.}
Пусть два тела $A$ и $B$ движутся прямолинейно и равномерно.
В системе отсчёта $A$ тело $B$ движется с постоянной скоростью
$\mbfv_{BA} = \mbfv_B - \mbfv_A$ по прямой линии.
Минимальное расстояние между $A$ и $B$~--- это расстояние
от неподвижного $A$ до этой прямой, то есть длина перпендикуляра,
опущенного из $A$ на прямую траектории $B$ в системе $A$.

\begin{example}
Два корабля движутся с постоянными скоростями $v_1 = v_2 = v$.
В некоторый момент расстояние между кораблями равно $L$;
их взаимное расположение и скорости показаны на рисунке:
корабль $A$ движется вертикально вверх, корабль $B$~---
под углом $30^\circ$ к горизонтали, причём $B$ находится строго
справа от $A$ на расстоянии $L$. Определите минимальное
расстояние между кораблями.

Переходим в систему отсчёта $A$. Скорость $B$ относительно $A$:
\[
    \mbfv_{BA} = \mbfv_B - \mbfv_A.
\]
Вектор $\mbfv_A$ направлен вертикально вверх, $|\mbfv_A| = v$.
Вектор $\mbfv_B$ направлен под $30^\circ$, $|\mbfv_B| = v$.
Разложим по компонентам (ось $x$~--- горизонталь, ось $y$~--- вертикаль):
$\mbfv_A = \colvecii{0}{v}$, $\mbfv_B = \colvecii{v\cos 30^\circ}{v\sin 30^\circ} = \colvecii{\dfrac{v\sqrt{3}}{2}}{\dfrac{v}{2}}$.
Тогда относительная скорость:
\[
    \mbfv_{BA} = \colvecii{\frac{v\sqrt{3}}{2}}{-\frac{v}{2}},
    \quad |\mbfv_{BA}| = v.
\]
В системе $A$ корабль $B$ стартует из точки $(L, 0)$
и движется в направлении единичного вектора
$\hat{\mbfn} = \colvecii{\dfrac{\sqrt{3}}{2}}{-\dfrac{1}{2}}$.
Минимальное расстояние~--- перпендикуляр из начала координат
на прямую движения $B$. В двумерном случае оно вычисляется как
модуль векторного произведения начального радиус-вектора
и направляющего вектора:
\[
    d_{\min} = \left| \mbfr_0 \times \hat{\mbfn} \right|
             = \left| L \cdot \left(-\frac{1}{2}\right) - 0 \cdot \frac{\sqrt{3}}{2} \right|
             = \frac{L}{2}.
\]
Минимальное расстояние между кораблями равно $\frac{L}{2}$.
\end{example}

\begin{example}
Пловец хочет переплыть реку шириной $L$.
Скорость течения равна $u$, скорость пловца относительно воды~--- $v$ ($v > u$).
Под каким углом к берегу он должен плыть, чтобы переправиться
за наименьшее время? Какое расстояние он пройдёт?
За какое время пловец переплывёт реку по наикратчайшему пути?

\noindent\textit{Наименьшее время.}
Время переправы~--- это ширина реки, делённая на поперечную
компоненту абсолютной скорости пловца.
Абсолютная скорость: $\vec{v}_{\text{абс}} = \vec{v}_{\text{отн}} + \vec{u}$,
где $\vec{u}$~--- скорость течения (горизонтально вдоль берега).
Поперечная компонента $\vec{v}_{\text{абс}}$ равна поперечной компоненте
$\vec{v}_{\text{отн}}$ (так как $\vec{u}$ направлено вдоль берега).
Она максимальна, если пловец плывёт строго перпендикулярно
берегу ($\alpha = 90^\circ$ к направлению течения):
$v_\perp^{\max} = v$, время $t_{\min} = \frac{L}{v}$.

При этом течение сносит пловца на $S_x = u \cdot t_{\min} = \frac{uL}{v}$.
Пройденное расстояние: 
\[
    S = \sqrt{L^2 + \left(\frac{uL}{v}\right)^2} = \frac{L}{v}\sqrt{u^2+v^2}.
\]

\noindent\textit{Кратчайший путь.}
Для прохождения минимального пути нужно так выбрать направление
$\vec{v}_{\text{отн}}$, чтобы абсолютная скорость была направлена
строго поперёк. Из треугольника скоростей: $\vec{v}_{\text{абс}} \perp \vec{u}$,
если $\vec{v}_{\text{отн}}$ направлен под углом $\sin\varphi = \frac{u}{v}$
против течения. Тогда $v_\perp = v\cos\varphi = \sqrt{v^2 - u^2}$, время:
\[
    t = \frac{L}{\sqrt{v^2 - u^2}}.
\]
\end{example}

%==============================================================================
\subsection{Упругое отражение от движущейся поверхности}
%==============================================================================

Когда мяч упруго ударяется о неподвижную гладкую стенку, нормальная
компонента скорости меняет знак, тангенциальная остаётся прежней.
Если стенка движется, задача в лабораторной системе выглядит сложнее.
Однако переход в систему стенки сводит её к уже известной задаче.

\noindent\textbf{Метод.}
Пусть стенка движется со скоростью $\vec{u}$ (перпендикулярно или
под углом). Перейдём в систему отсчёта стенки~--- там она покоится.
В этой системе мяч имеет скорость $\vec{v}' = \vec{v} - \vec{u}$.
После упругого удара о неподвижную стенку нормальная компонента
$\vec{v}'_n$ меняет знак; тангенциальная $\vec{v}'_\tau$ сохраняется.
Получив скорость $\vec{v}''$ в системе стенки, возвращаемся в
лабораторную: $\vec{v}_{\text{после}} = \vec{v}'' + \vec{u}$.

\begin{example}
Мяч движется перпендикулярно стенке со скоростью $v$;
стенка движется перпендикулярно самой себе навстречу мячу
со скоростью $u$ ($u < v$).
Найдите скорость мяча после удара.

Введём ось $x$, направленную от мяча к стенке.
Скорость мяча в лабораторной системе: $v_m = +v$.
Скорость стенки: $v_s = -u$ (движется навстречу).
Скорость мяча в системе стенки: 
\[
    v' = v_m - v_s = v - (-u) = v + u.
\]
После отражения (нормальная компонента меняет знак):
$v'' = -(v + u)$.
Скорость в лабораторной системе:
\[
    v_{\text{после}} = v'' + v_s = -(v+u) + (-u) = -v - 2u.
\]
Модуль скорости: $|v_{\text{после}}| = v + 2u$.
Мяч отлетает со скоростью $v + 2u$ в направлении, откуда прилетел.
Если стенка и мяч движутся в одну сторону ($v_s = +u$), аналогично
получим $|v_{\text{после}}| = |v - 2u|$.
\end{example}

\begin{example}
Упругий мячик летит в направлении движущейся от него стенки
со скоростью $v$. Стенка движется от мяча со скоростью $u$,
составляющей угол $\alpha = 30^\circ$ с нормалью к стенке; скорость
мяча вдвое больше скорости стенки: $v = 2u$.
Определите скорость мяча после удара.

Введём систему координат: ось $n$ перпендикулярна стенке
(от стенки к мячу), ось $\tau$~--- вдоль стенки.
Скорость стенки: $u_n = u\cos\alpha = \frac{u\sqrt{3}}{2}$~--- от мяча;
$u_\tau = u\sin\alpha = \frac{u}{2}$~--- вдоль стенки.

Скорость мяча в лабораторной системе направлена вдоль движения стенки,
то есть $(v_n, v_\tau) = (-v\cos\alpha,\; -v\sin\alpha) = (-u\sqrt{3},\; -u)$
(мяч летит в сторону стенки, поэтому компонента по $n$ отрицательна).

Скорость мяча в системе стенки:
\[
    (v'_n, v'_\tau) = (v_n - u_n,\; v_\tau - u_\tau) 
                    = \left(-u\sqrt{3} - \frac{u\sqrt{3}}{2},\; -u - \frac{u}{2}\right) 
                    = \left(-\frac{3u\sqrt{3}}{2},\; -\frac{3u}{2}\right).
\]
После отражения: $v''_n = -v'_n = \frac{3u\sqrt{3}}{2}$ (нормальная меняет знак),
$v''_\tau = v'_\tau = -\frac{3u}{2}$ (тангенциальная не меняется).

Возвращаемся в лабораторную систему:
\[
    v_{\text{после},n} = v''_n + u_n = \frac{3u\sqrt{3}}{2} + \frac{u\sqrt{3}}{2} = 2u\sqrt{3},
\]
\[
    v_{\text{после},\tau} = v''_\tau + u_\tau = -\frac{3u}{2} + \frac{u}{2} = -u.
\]
Итоговая скорость: $\vec{v}_{\text{после}} = (2u\sqrt{3},\; -u)$.
Её модуль: $|\vec{v}_{\text{после}}| = u\sqrt{12 + 1} = u\sqrt{13}$.
\end{example}

%==============================================================================
\subsection{Эффект Доплера}
%==============================================================================

Если источник звука движется, наблюдатель воспринимает иную частоту,
нежели та, с которой источник испускает звук.
Это явление называется \textit{эффектом Доплера}.
Оно не является чисто кинематическим, но сводится к нему,
поскольку скорость звука постоянна в данной среде
и не зависит ни от скорости источника, ни от скорости наблюдателя.

\noindent\textbf{Движущийся источник, неподвижный наблюдатель.}
Источник испускает звук с частотой $f_0$ (то есть испускает
$f_0$ волновых фронтов в секунду). Пусть источник движется
со скоростью $v_s$ в направлении наблюдателя, скорость звука~--- $c$.

За период $T_0 = \frac{1}{f_0}$ источник испускает один фронт
и успевает приблизиться к наблюдателю на $v_s T_0$.
Поэтому расстояние между соседними фронтами (длина волны
в лабораторной системе в направлении наблюдателя):
\[
    \lambda = c T_0 - v_s T_0 = (c - v_s) T_0 = \frac{c - v_s}{f_0}.
\]
Наблюдатель видит фронты, приходящие со скоростью $c$,
поэтому воспринимаемая частота:
\[
    f = \frac{c}{\lambda} = \frac{c \cdot f_0}{c - v_s}.
\]
Если источник удаляется ($v_s < 0$): $f = \frac{c f_0}{c + |v_s|}$~--- частота понижается.

\noindent\textbf{Неподвижный источник, движущийся наблюдатель.}
Теперь источник покоится; длина волны постоянна: $\lambda = \frac{c}{f_0}$.
Наблюдатель движется навстречу источнику со скоростью $v_r$.
Относительно наблюдателя фронты набегают со скоростью $c + v_r$,
поэтому воспринимаемая частота:
\[
    f = \frac{c + v_r}{\lambda} = \frac{(c + v_r) f_0}{c}.
\]

\noindent\textbf{Физическая причина асимметрии.}
Хотя в обоих случаях наблюдатель и источник сближаются с одной и
той же относительной скоростью, формулы различны.
Это объясняется тем, что звук~--- волна в среде (воздухе),
и его скорость $c$ задана относительно среды, а не
относительно источника или наблюдателя.
При движении источника изменяется \textit{длина волны}
(расстояние между фронтами); при движении наблюдателя
длина волны остаётся прежней, но наблюдатель встречает
фронты с иной скоростью.

\begin{definition}
Общая формула эффекта Доплера при движении в одной прямой
($v_s > 0$~--- источник движется к наблюдателю,
$v_r > 0$~--- наблюдатель движется к источнику):
\[
    f = f_0 \cdot \frac{c + v_r}{c - v_s}.
\]
\end{definition}

\begin{example}
С подводной лодки, погружающейся вертикально равномерно,
испускаются звуковые импульсы длительностью $\tau_0$.
Длительность приёма отражённого от дна импульса равна $\tau$;
скорость звука в воде $c$. Найдите скорость погружения лодки.

Рассмотрим импульс как серию волновых фронтов,
распространяющихся вниз. За время $\tau_0$ лодка опускается
на $v\tau_0$ (где $v$~--- искомая скорость).
Это задача на двукратный эффект Доплера.
Первый проход (вниз~--- источник движется к отражателю):
дно воспринимает частоту $f_1 = f_0 \cdot \frac{c}{c - v}$.
Второй проход (вверх~--- источник теперь дно, наблюдатель~--- лодка,
приближающаяся к отражённому фронту):
$f_2 = f_1 \cdot \frac{c + v}{c} = f_0 \cdot \frac{c + v}{c - v}$.

Поскольку импульс длительностью $\tau_0$ воспринимается
как импульс длительностью $\tau = \tau_0 \cdot \frac{f_0}{f_2} = \tau_0 \cdot \frac{c - v}{c + v}$,
(длительность обратно пропорциональна частоте) получаем:
\[
    \frac{\tau}{\tau_0} = \frac{c - v}{c + v}
    \quad\Longrightarrow\quad
    v = c \cdot \frac{\tau_0 - \tau}{\tau_0 + \tau}.
\]
Если лодка всплывает ($v < 0$ при обоих проходах), знаки
в числителе и знаменателе меняются местами:
$v = c \cdot \frac{\tau - \tau_0}{\tau + \tau_0}$.
\end{example}

\begin{example}
Наблюдатель сидит на крыше; рядом с ним~--- слушатель.
Вдали неподвижный дятел долбит дерево ровно $60$ ударов в минуту;
слушатель слышит $59$ ударов в минуту.
Скорость звука $c = 330$ м/с. Как движется крыша?

Дятел~--- источник с частотой $f_0 = 60$ мин$^{-1}$;
слушатель~--- движущийся наблюдатель с воспринимаемой
частотой $f = 59$ мин$^{-1}$.
Наблюдатель движется вместе с крышей. Из формулы:
\[
    f = f_0 \cdot \frac{c + v_r}{c}
    \quad\Longrightarrow\quad
    v_r = c\left(\frac{f}{f_0} - 1\right) = 330 \cdot \frac{59 - 60}{60} = -5{,}5~\text{м/с}.
\]
Знак минус означает, что крыша движется \textit{от} дятла.
\end{example}

%==============================================================================
\subsection{Вращающаяся система отсчёта: скорости}
%==============================================================================

До сих пор мы рассматривали системы отсчёта, движущиеся поступательно~---
их оси оставались параллельными лабораторным. Теперь рассмотрим
\textit{вращающуюся} систему~--- ту, чьи оси поворачиваются с течением времени.
Это принципиально важный шаг: именно вращающиеся системы описывают
наблюдателей, находящихся на Земле, на карусели или на вращающемся
диске.

\noindent\textbf{Вращающийся базис.}
Пусть система $S'$ вращается вокруг неподвижного начала координат
с постоянной угловой скоростью $\omega$ (для простоты рассматриваем плоское
движение). Её оси задаются единичными векторами $\hat{e}_1(t)$
и $\hat{e}_2(t)$, которые поворачиваются вместе
с системой.
В начальный момент $\hat{e}_1 = (1, 0)$, $\hat{e}_2 = (0, 1)$;
через время $t$:
\[
    \hat{e}_1(t) = (\cos\omega t,\; \sin\omega t),
    \qquad
    \hat{e}_2(t) = (-\sin\omega t,\; \cos\omega t).
\]
Вычислим производные по времени:
\[
    \dv{\hat{e}_1}{t} = \omega(-\sin\omega t,\; \cos\omega t) = \omega\,\hat{e}_2,
    \qquad
    \dv{\hat{e}_2}{t} = \omega(-\cos\omega t,\; -\sin\omega t) = -\omega\,\hat{e}_1.
\]
Это ключевой факт: \textit{производная единичного вектора, участвующего
во вращении, перпендикулярна самому вектору и по модулю равна $\omega$}.
В трёхмерном случае это записывается как $\dv{\hat{e}_i}{t} = \vec{\Omega}\times\hat{e}_i$,
где вектор угловой скорости $\vec{\Omega}$ перпендикулярен плоскости вращения.

\noindent\textbf{Вывод формулы для скоростей.}
Пусть точка $P$ имеет в системе $S'$ координаты $(x', y')$;
её радиус-вектор в лабораторной системе:
\[
    \vec{r} = x'\hat{e}_1 + y'\hat{e}_2.
\]
Дифференцируем по времени, используя правило Лейбница:
\[
    \vec{v} = \dv{\vec{r}}{t} = \dv{x'}{t}\hat{e}_1 + \dv{y'}{t}\hat{e}_2
    + x'\dv{\hat{e}_1}{t} + y'\dv{\hat{e}_2}{t}.
\]
Подставляем производные базисных векторов:
\[
    \vec{v} = \underbrace{\dv{x'}{t}\hat{e}_1 + \dv{y'}{t}\hat{e}_2}_{\vec{v}'} +
    x'\omega\hat{e}_2 - y'\omega\hat{e}_1
    = \vec{v}' + \omega(-y'\hat{e}_1 + x'\hat{e}_2).
\]
Выражение $\omega(-y'\hat{e}_1 + x'\hat{e}_2)$~--- это в точности
$\vec{\Omega}\times\vec{r}'$, где $\vec{\Omega} = \omega\hat{k}$
($\hat{k}$~--- орт, перпендикулярный плоскости).
Действительно, $\vec{\Omega}\times(x'\hat{e}_1 + y'\hat{e}_2)
= \omega\hat{k}\times(x'\hat{e}_1 + y'\hat{e}_2)
= \omega(x'\hat{e}_2 - y'\hat{e}_1)$~--- совпадает.

\begin{definition}
Закон сложения скоростей во вращающейся системе отсчёта
(начало координат неподвижно):
\[
    \vec{v} = \vec{v}' + \vec{\Omega}\times\vec{r}',
\]
где $\vec{v}'$~--- скорость точки в $S'$ (относительная),
$\vec{\Omega}\times\vec{r}'$~--- переносная скорость, обусловленная
вращением системы.
В общем случае, когда начало координат $O'$ движется со скоростью $\vec{V}_0$:
\[
    \vec{v} = \vec{V}_0 + \vec{v}' + \vec{\Omega}\times\vec{r}'.
\]
\end{definition}

Смысл этой формулы таков: скорость точки в лабораторной системе
складывается из трёх вкладов~--- скорости начала координат,
скорости точки в подвижной системе и скорости, порождённой
вращением подвижной системы.

\noindent\textbf{Связь с кинематикой твёрдого тела.}
Если точка $P$ покоится в системе $S'$ ($\vec{v}' = 0$),
то $\vec{v} = \vec{V}_0 + \vec{\Omega}\times\vec{r}'$~---
это в точности формула скоростей абсолютно твёрдого тела,
которую мы выводили в лекции о плоскопараллельном движении.
Там $\vec{\Omega}$ играло роль угловой скорости тела.
Мы видим, что поле скоростей твёрдого тела~---
это частный случай переноса из вращающейся системы отсчёта,
связанной с телом, в лабораторную.

\begin{example}
Карусель вращается с угловой скоростью $\Omega$.
Ребёнок сидит на карусели на расстоянии $r$ от оси и бросает
шарик радиально от оси со скоростью $u$ (в системе карусели).
Найдите скорость шарика в лабораторной системе.

В системе карусели шарик движется со скоростью $\vec{v}' = u\hat{r}'$,
где $\hat{r}'$~--- единичный вектор от оси. Начало координат
на оси ($\vec{V}_0 = 0$). По формуле:
\[
    \vec{v} = u\hat{r}' + \vec{\Omega}\times r\hat{r}'.
\]
Вектор $\vec{\Omega}\times r\hat{r}'$ перпендикулярен $\hat{r}'$
и имеет модуль $\Omega r$ (скорость, обусловленная вращением
точки начального положения).
Итого: шарик летит по касательной с компонентами $u$ (радиально)
и $\Omega r$ (по касательной):
$|\vec{v}| = \sqrt{u^2 + \Omega^2 r^2}$.
\end{example}

\begin{example}
В предыдущем примере опишите траекторию шарика в системе карусели
после броска (без сил трения; шарик свободно скользит по карусели).

В лабораторной системе на шарик (после броска) не действуют
горизонтальные силы~--- он движется прямолинейно. В системе
карусели прямая линия лабораторной системы выглядит кривой,
потому что наблюдатель вращается.
Именно это изогнутое движение~--- следствие \textit{сил инерции},
возникающих при описании движения во вращающейся системе.
Подробнее мы разберём это в следующем параграфе.
\end{example}

%==============================================================================
\subsection{Вращающаяся система отсчёта: ускорения. Ускорение Кориолиса}
%==============================================================================

Найдём, как связаны ускорения точки в лабораторной и вращающейся
системах отсчёта. Для этого продифференцируем формулу для скоростей.

\noindent\textbf{Ключевое правило дифференцирования.}
Прежде чем дифференцировать, сформулируем важное правило.
Если $\vec{A}(t)$~--- произвольный вектор, заданный своими
компонентами в вращающейся системе ($\vec{A} = A_1'\hat{e}_1 + A_2'\hat{e}_2$),
то его производная в лабораторной системе:
\[
    \dv{\vec{A}}{t} = \dv{A_1'}{t}\hat{e}_1 + \dv{A_2'}{t}\hat{e}_2
    + A_1'\dv{\hat{e}_1}{t} + A_2'\dv{\hat{e}_2}{t}
    = \left.\dv{\vec{A}}{t}\right|_{S'} + \vec{\Omega}\times\vec{A},
\]
где $\left.\dv{\vec{A}}{t}\right|_{S'} = \dv{A_1'}{t}\hat{e}_1 + \dv{A_2'}{t}\hat{e}_2$~---
производная вектора в системе $S'$ (только компоненты меняются,
базис фиксирован).
Это правило справедливо для \textit{любого} вектора, не только
для радиус-вектора, и является фундаментом теории вращающихся
систем отсчёта.

\noindent\textbf{Вывод формулы для ускорений.}
Начинаем с формулы для скоростей ($\vec{V}_0 = 0$ для простоты):
\[
    \vec{v} = \vec{v}' + \vec{\Omega}\times\vec{r}'.
\]
Дифференцируем по времени, применяя ключевое правило к каждому слагаемому.

Для первого слагаемого $\vec{v}'$:
\[
    \dv{\vec{v}'}{t} = \left.\dv{\vec{v}'}{t}\right|_{S'} + \vec{\Omega}\times\vec{v}' = \vec{a}' + \vec{\Omega}\times\vec{v}'.
\]

Для второго слагаемого $\vec{\Omega}\times\vec{r}'$ ($\vec{\Omega} = \mathrm{const}$):
\[
    \dv{(\vec{\Omega}\times\vec{r}')}{t} = \vec{\Omega}\times\dv{\vec{r}'}{t}.
\]
Производная $\vec{r}'$ в лабораторной системе~--- это $\vec{v}$ (полная скорость).
Но нам нужна производная в $S'$, чтобы применить правило: используем,
что $\dv{\vec{r}'}{t} = \vec{v}' + \vec{\Omega}\times\vec{r}'$
(из предыдущего параграфа). Тогда:
\[
    \dv{(\vec{\Omega}\times\vec{r}')}{t} = \vec{\Omega}\times(\vec{v}' + \vec{\Omega}\times\vec{r}') = \vec{\Omega}\times\vec{v}' + \vec{\Omega}\times(\vec{\Omega}\times\vec{r}').
\]

Складываем:
\[
    \vec{a} = \vec{a}' + \vec{\Omega}\times\vec{v}' + \vec{\Omega}\times\vec{v}' + \vec{\Omega}\times(\vec{\Omega}\times\vec{r}').
\]

\begin{definition}
Формула Кориолиса для ускорений (начало координат неподвижно,
$\vec{\Omega} = \mathrm{const}$):
\[
    \vec{a} = \vec{a}' + 2\vec{\Omega}\times\vec{v}' + \vec{\Omega}\times(\vec{\Omega}\times\vec{r}').
\]
Слагаемые справа~--- относительное ускорение $\vec{a}'$,
\textit{кориолисово ускорение} $\vec{a}_\kappa = 2\vec{\Omega}\times\vec{v}'$
и \textit{центростремительное переносное ускорение}
$\vec{a}_{\text{цс}} = \vec{\Omega}\times(\vec{\Omega}\times\vec{r}')$.
\end{definition}

\noindent\textbf{Физический смысл слагаемых.}
Центростремительное ускорение $\vec{\Omega}\times(\vec{\Omega}\times\vec{r}')$~---
это ускорение, которое имела бы точка, если бы покоилась в системе $S'$.
Оно направлено к оси вращения (перпендикулярно $\vec{\Omega}$)
и по величине равно $-\Omega^2 \vec{r}'_\perp$, где $\vec{r}'_\perp$~---
компонента $\vec{r}'$, перпендикулярная $\vec{\Omega}$
(расстояние от оси вращения).

Кориолисово ускорение $2\vec{\Omega}\times\vec{v}'$ возникает только
при наличии относительного движения ($\vec{v}' \ne 0$).
Оно перпендикулярно как $\vec{\Omega}$, так и $\vec{v}'$,
и по модулю равно $2\Omega v'\sin\theta$, где $\theta$~---
угол между $\vec{\Omega}$ и $\vec{v}'$.
Для наблюдателя, находящегося во вращающейся системе,
оно проявляется как сила ($\vec{F}_\kappa = -m\vec{a}_\kappa = -2m\vec{\Omega}\times\vec{v}'$)~---
\textit{сила Кориолиса}.

\begin{example}
Почему ураганы в Северном полушарии вращаются против часовой стрелки?

Воздух втекает в область пониженного давления (центр урагана).
В системе отсчёта, связанной с Землёй, он движется с ненулевой
относительной скоростью $\vec{v}'$, и на него действует сила Кориолиса.
Земля вращается с запада на восток; вектор $\vec{\Omega}$
направлен на Северный полюс (из центра Земли).
Для воздуха, движущегося к югу, сила Кориолиса:
$\vec{F}_\kappa = -2m\vec{\Omega}\times\vec{v}'$.
Вычисляем направление: $\vec{\Omega}$ смотрит вверх-на-север,
$\vec{v}'$ смотрит на юг; $\vec{\Omega}\times\vec{v}'$~--- на запад;
$\vec{F}_\kappa = -2m\cdot(\text{на запад}) = \text{на восток}$.
Для воздуха, движущегося с севера, он отклоняется на восток,
формируя вращение против часовой стрелки вокруг центра.
В Южном полушарии $\vec{\Omega}$ меняет направление (относительно
локальной вертикали), и вращение обращается по часовой стрелке.
\end{example}

\begin{example}
Тело падает с высоты $h$ без начальной скорости.
Найдите отклонение точки падения от точки прямо под местом
брошения (пренебрегите атмосферным сопротивлением,
учтите вращение Земли, широта~--- $\varphi$).

В системе отсчёта Земли на тело действует сила Кориолиса.
В начале падения $\vec{v}' \approx 0$, поэтому кориолисово ускорение
мало и нарастает по мере набора скорости. В первом приближении
вертикальная скорость: $v_z' \approx gt$ (вниз).

Горизонтальная составляющая $\vec{\Omega}$ на широте $\varphi$,
направленная на север: $\Omega_h = \Omega\cos\varphi$.
Кориолисово ускорение (горизонтальное, на восток):
$a_\kappa = 2\Omega\cos\varphi \cdot v_z' = 2\Omega\cos\varphi \cdot gt$.

Горизонтальная скорость на восток (нарастает):
\[
    \dv{x'}{t} \approx \int_0^t 2\Omega\cos\varphi\cdot g t'\,\dd{t'} = \Omega g\cos\varphi\cdot t^2.
\]

Горизонтальное смещение:
\[
    x' = \int_0^T \Omega g\cos\varphi\cdot t^2\,\dd{t} = \frac{\Omega g\cos\varphi}{3} T^3.
\]

Время падения $T = \sqrt{\frac{2h}{g}}$, подставляем:
\[
    x' = \frac{\Omega\cos\varphi}{3}\sqrt{\frac{8h^3}{g}}.
\]
При $h = 100$ м, $\varphi = 45^\circ$, $\Omega = 7{,}3\cdot10^{-5}$ рад/с:
$x' \approx 16$ мм~--- небольшое, но измеримое отклонение на восток.
\end{example}

%==============================================================================
\subsection{Общий случай: теорема Кориолиса}
%==============================================================================

Соберём теперь все полученные результаты в единую формулу для
произвольной неинерциальной системы отсчёта~--- такой,
которая и движется поступательно, и вращается, и может иметь
переменную угловую скорость.

\noindent\textbf{Постановка задачи.}
Пусть начало координат $O'$ движется произвольно:
его скорость $\vec{V}_0(t)$, ускорение $\vec{a}_0(t) = \dv{\vec{V}_0}{t}$.
Одновременно система $S'$ вращается~--- угловая скорость
$\vec{\Omega}(t)$ может быть переменной.

\noindent\textbf{Формула скоростей.}
Из предыдущих параграфов, учитывая теперь и движение начала координат:
\[
    \vec{v} = \vec{V}_0 + \vec{v}' + \vec{\Omega}\times\vec{r}'.
\]

\noindent\textbf{Вывод формулы ускорений.}
Дифференцируем формулу скоростей:
\[
    \vec{a} = \vec{a}_0 + \dv{\vec{v}'}{t} + \dv{(\vec{\Omega}\times\vec{r}')}{t}.
\]
Первое слагаемое: $\vec{a}_0 = \dv{\vec{V}_0}{t}$~--- ускорение начала.

Второе: $\dv{\vec{v}'}{t} = \vec{a}' + \vec{\Omega}\times\vec{v}'$
(применяем ключевое правило к вектору $\vec{v}'$).

Третье: $\vec{\Omega}$ теперь переменно, поэтому:
\[
    \dv{(\vec{\Omega}\times\vec{r}')}{t} = \dv{\vec{\Omega}}{t}\times\vec{r}' + \vec{\Omega}\times\dv{\vec{r}'}{t}
    = \dv{\vec{\Omega}}{t}\times\vec{r}' + \vec{\Omega}\times(\vec{v}' + \vec{\Omega}\times\vec{r}').
\]

Складываем всё:

\begin{definition}
\textit{Теорема Кориолиса.}
Ускорение точки в лабораторной инерциальной системе:
\[
    \vec{a} = \vec{a}_0 + \vec{a}' + \underbrace{2\vec{\Omega}\times\vec{v}'}_{\text{Кориолис}} + \underbrace{\vec{\Omega}\times(\vec{\Omega}\times\vec{r}')}_{\text{центростремительное}} + \underbrace{\dv{\vec{\Omega}}{t}\times\vec{r}'}_{\text{Эйлера}}.
\]
Слагаемые:
$\vec{a}_0$~--- переносное ускорение начала координат;
$\vec{a}'$~--- относительное ускорение точки в $S'$;
$2\vec{\Omega}\times\vec{v}'$~--- \textit{ускорение Кориолиса};
$\vec{\Omega}\times(\vec{\Omega}\times\vec{r}')$~--- \textit{переносное центростремительное ускорение};
$\dv{\vec{\Omega}}{t}\times\vec{r}'$~--- \textit{ускорение Эйлера} (при переменной $\vec{\Omega}$).
\end{definition}

\noindent\textbf{Физический смысл.}
Переходя в систему $S'$ и переписывая уравнение $\vec{F} = m\vec{a}$
(записанное в инерциальной системе), мы получаем второй закон
Ньютона в неинерциальной системе:
\[
    m\vec{a}' = \vec{F} - m\vec{a}_0 - 2m\vec{\Omega}\times\vec{v}' - m\vec{\Omega}\times(\vec{\Omega}\times\vec{r}') - m\dv{\vec{\Omega}}{t}\times\vec{r}'.
\]
Наблюдатель в $S'$ видит дополнительные «силы»~--- это \textit{силы инерции}.
Они не являются настоящими силами (у них нет парных реакций),
но в расчётах ведут себя точно так же, как настоящие.

\begin{example}
Найдите эффективное ускорение свободного падения на поверхности
Земли как функцию широты $\varphi$.

Земля~--- наша система $S'$; она вращается с $\vec{\Omega}$
и сама движется вокруг Солнца (этим движением пренебрежём:
ускорение орбитального движения $a_0 \approx 6\cdot10^{-3}$~м/с$^2$,
что значительно меньше $g$).
Рассмотрим покоящееся тело ($\vec{v}' = 0$, $\vec{a}' = \vec{g}_{\text{eff}}$).
Из теоремы Кориолиса (кориолисово и эйлерово ускорения исчезают):
\[
    \vec{g}_{\text{eff}} = \vec{g}_{\text{ист}} - \vec{\Omega}\times(\vec{\Omega}\times\vec{R}),
\]
где $\vec{g}_{\text{ист}} = -\frac{GM}{R^2}\hat{r}$~--- истинное
гравитационное ускорение, $\vec{R}$~--- радиус-вектор точки.

Вектор $\vec{\Omega}\times(\vec{\Omega}\times\vec{R})$ направлен
\textit{к оси вращения Земли} и имеет модуль $\Omega^2 R\cos\varphi$
(расстояние до оси умножить на $\Omega^2$).
На экваторе $\varphi = 0$, это ускорение максимально и направлено
от центра Земли: $\Omega^2 R \approx 0{,}034$~м/с$^2$.
Оно уменьшает эффективное $g$:
$g_{\text{экватор}} = g_{\text{ист}} - \Omega^2 R \approx 9{,}78$~м/с$^2$.

На полюсе $\varphi = 90^\circ$~--- расстояние до оси равно нулю,
центростремительное слагаемое исчезает:
$g_{\text{полюс}} = g_{\text{ист}} \approx 9{,}83$~м/с$^2$.
Разница в $g$ между экватором и полюсом составляет около $0{,}05$~м/с$^2$~---
небольшая, но вполне измеримая величина.
Кроме того, на экваторе Земля из-за центробежного эффекта немного
сплюснута, что дополнительно уменьшает $g$ (бо́льшее расстояние
от центра уменьшает гравитационное притяжение).
\end{example}

\begin{example}
Маятник Фуко~--- длинный маятник, подвешенный в точке на широте $\varphi$.
Покажите, что плоскость его колебаний медленно вращается.

В системе отсчёта Земли на маятник действует сила Кориолиса
$\vec{F}_\kappa = -2m\vec{\Omega}\times\vec{v}'$.
Вертикальная составляющая $\vec{\Omega}$: $\Omega_v = \Omega\sin\varphi$.
Горизонтальная скорость маятника $\vec{v}'_h$ вызывает горизонтальную
силу Кориолиса, перпендикулярную этой скорости:
$\vec{F}_{\kappa,h} = -2m\Omega\sin\varphi\;\hat{k}\times\vec{v}'_h$.
Это~--- та же структура, что у силы Лоренца ($\vec{F} = q\vec{v}\times\vec{B}$),
поэтому маятник ведёт себя как заряженный осциллятор в магнитном поле:
его плоскость колебаний медленно прецессирует.

Угловая скорость прецессии~--- это угловая скорость вращения
$\Omega_v = \Omega\sin\varphi$. Период полного оборота плоскости:
\[
    T_{\text{Фуко}} = \frac{2\pi}{\Omega\sin\varphi} = \frac{24\text{ ч}}{\sin\varphi}.
\]
На Северном полюсе ($\varphi = 90^\circ$): $T = 24$ ч~--- плоскость маятника
остаётся неподвижной, пока Земля оборачивается под ним.
На экваторе ($\varphi = 0^\circ$): $T = \infty$~--- плоскость вообще не вращается.
В Санкт-Петербурге ($\varphi \approx 60^\circ$): $T \approx 27{,}7$ ч.
\end{example}

\end{document}